\SourceLeaf{203}
\[
\begin{aligned}
P&=\tfrac{1}{2}f(p+iq)+\tfrac{1}{2}f'(p-iq)\\
iQ&=\tfrac{1}{2}f(p+iq)-\tfrac{1}{2}f'(p-iq)
\end{aligned}
\]
\noindent oder, was dasselbe ist, indem die Function \(f\) ganz willkürlich angenommen wird (nach Gefallen mit Inbegriff constanter imaginärer Elemente), wird \(P\) dem reellen und \(iQ\) (bei der zweiten Auflösung \(-iQ\)) dem imaginären Theile von \(f(p+iq)\) gleich gesetzt, und hieraus sodann vermittelst der Elimination \(T\) und \(U\) in der Gestalt von Functionen von \(t\) und \(u\) dargestellt werden. Hiedurch ist die vorgegebene Aufgabe ganz allgemein und vollständig aufgelöst.

\Article{6}
Wenn \(p'+iq'\) eine beliebige bestimmte Function von \(p+iq\) vorstellt (indem \(p',q'\) reelle Functionen von \(p,q\) sind), so sieht man leicht, dass auch
\[
p'+iq'=\text{Const.}\quad\text{und}\quad p'-iq'=\text{Const.}
\]
\noindent die Integrale der Differentialgleichung \(\omega=0\) darstellen; in der That werden jene mit den obigen
\[
p+iq=\text{Const.}\quad\text{und}\quad p-iq=\text{Const.}
\]
\noindent resp. ganz gleichbedeutend sein. Eben so werden die Integrale der Differentialgleichung \(\Omega=0\)
\[
P'+iQ'=\text{Const.}\quad\text{und}\quad P'-iQ'=\text{Const.}
\]
\noindent mit den obigen
\[
P+iQ=\text{Const.}\quad\text{und}\quad P-iQ=\text{Const.}
\]
\noindent ganz gleichbedeutend sein, wenn \(P'+iQ'\) eine beliebige bestimmte Function von \(P+iQ\) vorstellt (indem \(P',Q'\) reelle Functionen von \(P,Q\) sind). Es erhellet hieraus, dass in der allgemeinen Auflösung unsrer Aufgabe, welche wir im vorhergehenden Artikel gegeben haben, auch \(p',q'\) die Stelle von \(p,q\); und \(P',Q'\) die Stelle von \(P,Q\) resp. vertreten können. Wenn gleich die Allgemeinheit der Auflösung durch eine solche Abänderung nichts gewinnt, so kann doch zuweilen für die Anwendung eine Form zu diesem, die andere zu jenem Zweck bequemer sein.
